\documentclass[11pt]{article}

\usepackage[margin=1in]{geometry}
\usepackage{times}
\usepackage{amsmath,amssymb}
\usepackage{xcolor}
\usepackage{enumitem}
\usepackage{microtype}
\usepackage[hidelinks]{hyperref}
\usepackage{parskip}

\setlength{\parindent}{0pt}
\setlist[itemize]{leftmargin=1.6em,itemsep=0.25em,topsep=0.35em}
\setlist[enumerate]{leftmargin=1.8em,itemsep=0.25em,topsep=0.35em}

\definecolor{reviewblue}{RGB}{0,74,135}

\newenvironment{reviewcomment}
  {\begin{quote}\color{reviewblue}\itshape}
  {\end{quote}}

\newcommand{\Response}{\par\noindent\textbf{Response.}\ }
\newcommand{\Changes}{\par\noindent\textbf{Changes in the revised manuscript.}\ }
\newcommand{\Location}[1]{\par\noindent\textit{Location in the revised manuscript: #1}}

\begin{document}

\begin{center}
{\LARGE\bfseries Cover Letter for Revised Submission}\\[0.8em]
{\large Manuscript ID: 26-3482}\\[0.35em]
{\large ``ADAPT: Adaptive Dual-projection Architecture for Perceptive Traversal''}
\end{center}

\vspace{0.8em}

Dear Editor, Associate Editor, and Reviewers,

Thank you for your careful assessment of our manuscript and for the constructive suggestions. We are revising the paper to clarify the comparison with voxel-based methods, the scope of adaptive sensing, and the presentation of the simulation and hardware experiments. The accompanying Statement of Changes provides point-by-point responses to the Associate Editor and both reviewers. Responses on deployment latency and onboard resource measurements, and on the causes of the voxel-based baselines' performance gap, will be completed after the corresponding experiments.

\end{document}
